\documentclass[11pt,a4paper]{article}
\usepackage[utf8]{inputenc}\usepackage[T1]{fontenc}\usepackage[polish]{babel}
\usepackage{mathptmx}
\usepackage[top=2.2cm,bottom=2.2cm,left=2.5cm,right=2.5cm]{geometry}
\usepackage{setspace}\onehalfspacing
\usepackage{parskip}\usepackage{microtype}
\usepackage{graphicx}\usepackage{float}\usepackage{booktabs}
\usepackage{tabularx}\usepackage{array}\usepackage{amssymb}\usepackage{amsmath}
\usepackage{xcolor}
\usepackage[hidelinks,colorlinks=true,linkcolor=blue!60!black,citecolor=blue!60!black]{hyperref}
\usepackage[numbers,sort&compress]{natbib}
\pagestyle{plain}

\title{\textbf{Czy modele językowe mogą symulować respondentów psychometrycznych?\\[4pt]
Badanie czułości na pięciu modelach z kontrolą bazową i kontrolą bez instrukcji\\dla polskiej baterii mentalizacyjnej}}
\author{Micha{\l} Wiencek\\\small Niezależny badacz; student studiów magisterskich, Instytut Psychologii UKEN, Krak\'ow}
\date{}

\begin{document}
\maketitle

\begin{abstract}
\noindent Zanim autor zbierze dane od ludzi do nowej baterii psychometrycznej, może najpierw zadać inne pytanie: czy narzędzia w ogóle reagują na zaprojektowane zróżnicowanie osobowościowe? Trzydzieści fikcyjnych biografii w języku polskim z ustalonymi profilami w 12 wymiarach zostało przetworzonych przez pięć dużych modeli językowych od czterech producentów --- Claude Sonnet 4.6 i Opus 4.6 (Anthropic), GPT-5.4-mini (OpenAI), Grok-4-20-reasoning (xAI), Gemini 3 Flash (Google) --- w trzech warunkach obserwacji: \textbf{warunek z personą} (każda biografia prezentowana dwukrotnie każdemu modelowi; 285 ważnych przebiegów), \textbf{warunek bazowy} (każdy model wypełniał baterię jako on sam, 7--10 razy; 46 przebiegów) oraz \textbf{warunek bez instrukcji} (bateria prezentowana bez prompta systemowego; 31 ważnych przebiegów). Siedmiu ludzkich respondentów wypełniło tę samą baterię (eksploracyjnie). Jest to badanie czułości, nie walidacja.

Warunek z personą pokazuje, że cztery opublikowane polskie narzędzia samoopisowe odzyskują wbudowane profile z wysoką spójnością (dopasowanie stylu przywiązania $\kappa = .73$ dla czterech modeli, $.33$ dla GPT, którego omawiamy jako outliera; mediana korelacji rangowo-wynikowej $r = .71$ do $.82$; spójność z literaturą $8$--$10$/$10$). Zgodność między producentami w $z$-skorach kwestionariuszowych jest wysoka (mediana $r$ w parach $= .78$--$.96$); zgodność między producentami w autorsko zakluczonym TCTM-22 jest wyraźnie niższa (ICC$(2,1) = .25$ dla wszystkich pięciu modeli, ale $.39$ po wykluczeniu GPT jako test wrażliwości), z ośmioma na 22 pozycjami wykazującymi statystycznie istotną niezgodność po korekcji Bonferroniego. Warunek bazowy pokazuje, że każdy model ma stabilną sygnaturę odpowiedzi w nieobecności persony (GPT klasyfikowany jako lękowo-ambiwalentny w 100\% przebiegów bazowych; Gemini osiąga sufit każdej skali samoopisowej; dwa modele Claude'a różnią się domyślnym stylem przywiązania) oraz że instrukcje persona rozszerzają wariancję odpowiedzi $3{,}5$ do $8{,}9$ razy w stosunku do wariancji bazowej (mediana na poziomie modelu w 9 wymiarach kwestionariuszowych). Warunek bez instrukcji pokazuje, że sama instrukcja ,,odpowiadaj jako sam siebie'' wiąże się z systematycznymi przesunięciami w samoopisie: względem braku instrukcji, instrukcja współwystępuje z około $0{,}9$ punktu skali niższym unikaniem przywiązaniowym dla Sonneta ($\Delta_{\text{Glassa}} \approx 1{,}96$) oraz $1{,}0$ punktu skali wyższą stabilnością emocjonalną dla Gemini (wewnątrzwarunkowe $\Delta = -1{,}00$); GPT odwraca kierunek na samoopisowej mentalizacji i lęku. Średnia trafność TCTM-22 zmienia się o mniej niż $\pm 5$ punktów procentowych między warunkiem bazowym a warunkiem bez instrukcji dla czterech z pięciu modeli; GPT spada o $8{,}1$ p.p.\ bez instrukcji.

Otwarte dane, biografie i skrypty analityczne są udostępnione na licencji CC BY 4.0; każde liczbowe twierdzenie w tej pracy jest reprodukowalne z wydanego pliku CSV.

\vspace{4pt}\noindent\textbf{Słowa kluczowe:} respondenci syntetyczni, czułość narzędzi, mentalizacja, przywiązanie, porównanie LLM między producentami, sygnatura bazowa odpowiedzi
\end{abstract}

% ═══════════════════════════════════
\section{Wstęp}

Tworzenie nowego narzędzia psychometrycznego wymaga rozległych badań pilotażowych. Rosnąca linia badań sugeruje, że duże modele językowe (LLM) mogą służyć jako symulowani respondenci, dostarczając szybkich pre-pilotowych sond zachowania narzędzi, zanim zaangażuję się w rekrutację ludzi \cite{argyle2023,dillion2023}. Wąskie pytanie tej pracy nie brzmi czy LLM ,,mają'' osobowość, lecz czy narzędzia psychometryczne \emph{reagują systematycznie} na zaprojektowane zróżnicowanie osobowościowe wbudowane w prompty biograficzne. Jeśli reagują, są przynajmniej czułe na odpowiedni sygnał --- warunek konieczny trafności konstruktu \cite{cronbach1955}. Jeśli nie, problem warto wychwycić przed rekrutacją ludzi.

Większość istniejących badań używa pojedynczej rodziny modeli \cite{argyle2023,dillion2023}; my używam pięciu modeli od czterech producentów (Anthropic, OpenAI, xAI, Google), co pozwala zapytać, czy ustalenia o czułości generalizują się między architekturami, czy są artefaktami specyficznymi dla modelu. Włączam też dwa warunki kontrolne: \textbf{warunek bazowy}, w którym model odpowiada na baterię jako on sam, wielokrotnie (sondując sygnaturę odpowiedzi modelu w nieobecności persony), oraz \textbf{warunek bez instrukcji}, w którym model otrzymuje baterię bez żadnej instrukcji systemowej (sondując, jaka część profilu odpowiedzi jest kształtowana nawet przez minimalny kontekst instrukcyjny). Razem porównanie persona / bazowy / bez instrukcji pozwala oddzielić trzy źródła wariancji odpowiedzi: nieodłączną stochastyczność wewnątrz warunku, wkład sformułowania instrukcyjnego oraz wkład promptu biograficznego.

Stosuję ten projekt do czterech opublikowanych polskich narzędzi --- DBZ-R (przywiązanie; polska adaptacja ECR-R), MentS-PL (mentalizacja), KPP (potrzeba poznania), TIPI-PL (Wielka Piątka) --- plus jednego oryginalnego, autorsko zakluczonego testu winietowego mentalizacji (TCTM-22). Trzydzieści fikcyjnych person jest przetwarzanych dwukrotnie przez każdy model; siedmiu ludzkich respondentów wypełnia tę samą baterię jako małe porównanie eksploracyjne.

% ═══════════════════════════════════
\section{Metoda}

\textbf{Persony.} 30 biografii w języku polskim (1{,}554--2{,}914 słów łącznie z metadanymi YAML; 1{,}489--2{,}861 słów narracyjnych po usunięciu YAML), napisanych w drugiej osobie, kodujących ustalone profile w 12 wymiarach (11 ciągłych + 1 kategoryczny): lęk i unikanie przywiązaniowe (DBZ-R), mentalizacja siebie / innych / motywacji (MentS-PL), potrzeba poznania (KPP), Wielka Piątka E/A/C/ES/O (TIPI-PL) oraz czterokategoryjny styl przywiązania. Biografie zostały napisane przez pierwszego autora przy pomocy Claude'a Opus 4.6 jako asystenta pisarskiego; autor specyfikował profile, ograniczenia narracyjne i wymagania stylistyczne, a także recenzował i edytował każdą wersję roboczą. Sześć biografii (\texttt{adrian, ewa, kuba, magda, pawel, weronika}) zostało zrewidowanych po wstępnej recenzji i jest oznaczonych \texttt{persona\_version: v2}; pozostałe są w wersji \texttt{v1}. Wersja wysłana do wszystkich pięciu modeli to wersja zarejestrowana w opublikowanym \texttt{all\_data.csv}. Żadna biografia nie była zaprojektowana tak, by zawierać etykiety cech; pierwszy autor ręcznie oznaczył dziewięć biografii jako zawierających bezpośrednie odniesienia do języka cech (\texttt{ania, anna-sim, filip, jakub, kamil, marek, natalia, ola, tomek}; podzbiór ,,explicit'') i 21 jako ,,narrative''. Osobny audyt regex z szerszym zestawem słów kluczowych daje inny podział; oba kodowania są w otwartym zbiorze danych.

\textbf{Narzędzia.} DBZ-R (polska adaptacja ECR-R; 36 pozycji, Likert 1--7; \citealt{lubiewska2016}); MentS-PL (28 pozycji; \citealt{janczak2021}); KPP (36 pozycji; \citealt{matusz2011}); TIPI-PL (10 pozycji; \citealt{sorokowska2014}). Te cztery narzędzia korzystają z opublikowanych norm polskich. TCTM-22 jest narzędziem oryginalnym: 22 winiety oceniane względem kluczy odpowiedzi zdefiniowanych przez autora, z każdą opcją oznaczoną według taksonomii błędów MASC \cite{dziobek2006}: \emph{DOS} (niedostateczna mentalizacja, redukowana teoria umysłu), \emph{NAD} (nadmierna mentalizacja, nadmiarowa teoria umysłu) oraz \emph{BK} (brak mentalizacji). 22 winiety obejmują pięć podskal (4--5 pozycji każda): rozpoznawanie podtekstu, rozpoznawanie manipulacji, wykrywanie sygnałów wygasania więzi, dynamika dążenie--wycofanie oraz ocena autentyczności naprawy relacyjnej. \textbf{W całej tej pracy wszystkie wyniki TCTM-22 odnoszą się do trafności względem klucza autorskiego, nie do zwalidowanego pomiaru mentalizacji}; niezgodność między producentami co do tego, które pozycje są trudne (\S\ref{sec:miedzyproducenci}), jest sama w sobie dowodem, że klucze odpowiedzi mogą nie odzwierciedlać pojedynczego wspólnego konstruktu.

\textbf{Prawda referencyjna i STYLE\_MAP.} Oczekiwany styl przywiązania każdej persony pochodzi z pola \texttt{expected\_style} na poziomie głównym w \texttt{analysis\_summary.jsonl}. Cztery persony (Ewa, Kamil, Marek, Radek) niosą \texttt{disorganized} w bloku YAML \texttt{expected\_profile}, ale \texttt{fearful\_avoidant} w polu \texttt{expected\_style}; mapuje to przywiązanie zdezorganizowane na bojaźliwo-unikowe w czterokategoryjnej taksonomii Bartholomew, której używają wszystkie pięć modeli w swoich klasyfikacyjnych wyjściach. Pod ścisłą polityką pięciokategoryjną (traktowanie \texttt{disorganized} jako piątej klasy, której żaden model nie potrafi wyprodukować) wszystkie wartości $\kappa$ przywiązania spadłyby o około 10 punktów procentowych; raportuję oba w opublikowanym archiwum suplementarnym.

\textbf{Modele.} Pięć modeli od czterech producentów (Tabela~\ref{tab:models}). Każda persona była przetwarzana raz przez każdy model (przebieg~1) i ponownie w identycznych warunkach (przebieg~2, retest). Wszystkie modele używały domyślnych parametrów próbkowania. Łącznie: 285 ważnych przebiegów persona z 300 prób; 15 brakujących retestów wynikało z filtrów moderacji treści API, częściowych generacji lub niekompletnych odpowiedzi TCTM (rozbicie per model w przypisie Tabeli~\ref{tab:models}). Pięć dodatkowych zapisów Gemini przebieg-3 (dla \texttt{ania, anna-sim, kasia, marek, michal-sim}) naprawia wcześniejsze zdegenerowane wyjścia i jest udostępnionych dla przejrzystości, ale są wykluczone ze wszystkich analiz przebieg~1 i przebieg~2, by uniknąć asymetrycznej nadreprezentacji. GPT-5.4-mini zwracał niekompletne odpowiedzi TCTM (20--21 z 22 pozycji) w 20\% prób; dotknięte przebiegi są wykluczone z tabel TCTM dla danego modelu ($N = 27$ przebieg~1, $N = 21$ przebieg~2). Dla czterech narzędzi samoopisowych brakujące pozycje są NA, z usuwaniem parami w analizach korelacji.

\begin{table}[htbp]\centering\footnotesize
\caption{Modele, producenci, platformy i liczby przebiegów po wykluczeniu częściowych prób / odfiltrowanych przez moderację.}
\label{tab:models}
\begin{tabular}{llll}
\toprule
\textbf{Model} & \textbf{Producent} & \textbf{Platforma} & \textbf{Przebiegi (prz.~1 + prz.~2)} \\
\midrule
Claude Sonnet 4.6 & Anthropic & AWS Bedrock & $30 + 30$ \\
Claude Opus 4.6 & Anthropic & AWS Bedrock & $30 + 25^{\ddagger}$ \\
GPT-5.4-mini & OpenAI & Azure OpenAI & $27^{\dagger} + 21^{\dagger}$ \\
Grok-4-20-reasoning & xAI & Azure AI Foundry & $30 + 29^{\ddagger}$ \\
Gemini 3 Flash & Google & Gemini API & $30 + 28^{\ddagger}$ ($+5$ prz.-3) \\
\midrule
Respondenci ludzcy & --- & --- & 7 \\
\bottomrule
\end{tabular}
\begin{flushleft}
\footnotesize $^{\dagger}$Filtr moderacji Azure lub częściowa odpowiedź TCTM. $^{\ddagger}$Jedna lub więcej person nie ukończyło przebiegu-2 z powodu przejściowych problemów API. Brakujące w przebiegu-2 Opus: \texttt{anna-sim, jakub, kasia, michal-sim, piotr}; brakujące w przebiegu-2 Grok: \texttt{piotr}; brakujące w przebiegu-2 Gemini: \texttt{ewa, sara}.
\end{flushleft}
\end{table}

\textbf{Warunki.} Uruchomiono trzy warunki.

\emph{Warunek z personą} (przebieg~1 + przebieg~2): model otrzymuje biografię w prompcie użytkownika oraz instrukcję systemową \textit{,,Odpowiadasz jako osoba opisana poniżej. Przeczytaj uważnie biografię i odpowiadaj tak, jak ta osoba by odpowiedziała.''}; 285 ważnych przebiegów.

\emph{Warunek bazowy}: każdy model wypełnia baterię 7--10 razy \emph{jako on sam}, z instrukcją systemową \textit{,,Jesteś modelem językowym AI. Odpowiadaj na pytania szczerze, jak Ty --- AI --- byś odpowiedział, gdybyś musiał wybrać. Nie odgrywaj żadnej postaci. Odpowiadaj jako sam siebie.''}; 46 ważnych przebiegów (jeden bazowy Gemini wykluczony za niekompletny TCTM).

\emph{Warunek bez instrukcji}: model otrzymuje tylko treść testu jako wiadomość użytkownika, bez prompta systemowego; 31 ważnych ładunków TCTM-22 (Sonnet 6, Opus 5, GPT 6, Grok 6, Gemini 8) i 30 parsowalnych ładunków samoopisowych (jedno wyjście GPT bez instrukcji zawierało kompletną tablicę TCTM-22, ale brak parsowalnych sekcji samoopisowych).

\textbf{Pilotaż na ludziach.} Siedmiu respondentów (wiek 22--60, 4 mężczyzn, 3 kobiety; 2 studentów psychologii, 5 bez wykształcenia psychologicznego) wypełniło tę samą baterię online. Docelowa populacja dla ewentualnej walidacji to młodzi dorośli (18--35) w romantycznych lub bliskich związkach; jeden respondent (R07, 60 lat, kobieta) był poza tym zakresem i jest zachowany wyłącznie dla celów eksploracyjnych. W całej pracy traktuję pilotaż na ludziach jako informację opisową, nie inferencyjną.

\textbf{Etyka.} Pilotaż na ludziach wykorzystał świadomą zgodę, ujawniającą cel badania, dobrowolność (prawo do wycofania w dowolnym momencie bez konsekwencji), zgodne z RODO przetwarzanie danych (szyfrowane przechowywanie, pięcioletni okres retencji, brak łączenia z osobistymi identyfikatorami) oraz linki do zasobów wsparcia zdrowia psychicznego (numery polskich krajowych telefonów kryzysowych). Wszystkich siedmiu respondentów wyraziło zgodę przed wypełnieniem baterii. Pilotaż jest nieinterwencyjny i nie używa populacji klinicznej. Autor nie konsultował formalnie z Komisją Bioetyczną UKEN statusu zwolnienia; czytelnicy powinni traktować część pilotażu na ludziach jako pre-pilotowe dane eksploracyjne w oczekiwaniu na taką konsultację.

\textbf{Mierniki.} (1) Dopasowanie stylu przywiązania ($\kappa$, $V$ Cramera); (2) korelacje rangowo-wynikowe (Pearson, Spearman); (3) spójność z literaturą (10 znaków korelacji między skalami pochodzących z opublikowanych norm polskich); (4) stabilność test--retest ($r$ Pearsona, ICC); (5) spójność między modelami (ICC, $\kappa$ Fleissa); (6) profile błędów TCTM-22 (DOS / NAD / BK); (7) zgodność na poziomie pozycji między producentami ($Q$ Cochrana z korekcjami Bonferroniego i Benjaminiego--Hochberga); (8) oszacowania wielkości efektu dla przesunięć sformułowania ($d$ Cohena pooled i $\Delta$ Glassa względem bazowego $SD$).

\textbf{Powtarzalność.} Wszystkie pięć modeli było odpytywanych przez oficjalne API producentów między 12 a 27 kwietnia 2026, z domyślnymi parametrami próbkowania (temperature, top-$p$, max-output-tokens przy wartościach domyślnych dostawców). Prompty systemowe i schematy JSON wyjścia są reprodukowane verbatim w otwartym repozytorium. Każde liczbowe twierdzenie w tej pracy jest reprodukowalne z opublikowanego \texttt{all\_data.csv} (369 obserwacji $\times$ 53 kolumny) używając wyłącznie standardowej biblioteki; skrypty pomocnicze są na licencji CC BY 4.0.

% ═══════════════════════════════════
\section{Wyniki}

Organizuję wyniki wokół trzech pytań. \textbf{Sekcja~\ref{sec:reaguja}}: czy narzędzia reagują na zaprojektowane zróżnicowanie? \textbf{Sekcja~\ref{sec:miedzyproducenci}}: czy reakcja generalizuje się między producentami? \textbf{Sekcja~\ref{sec:wariancja}}: skąd bierze się obserwowana wariancja odpowiedzi?

Wszystkie tabele między modelami w Sekcjach~\ref{sec:reaguja} i \ref{sec:miedzyproducenci} używają danych z pierwszego przebiegu warunku z personą ($N = 30$ na model z wyjątkiem GPT, gdzie trzy persony z niekompletnym TCTM są wykluczone, dając $N = 27$).

\subsection{Czy narzędzia reagują na zaprojektowane zróżnicowanie?}\label{sec:reaguja}

\subsubsection*{Rozkład wyników TCTM-22}

Pojawia się wyraźny gradient (Tabela~\ref{tab:tctm}, Rycina~\ref{fig:tctm}): pilotaż na ludziach uzyskał $M = 14{,}3$ ($SD = 1{,}4$), GPT $M = 17{,}0$, Grok $M = 18{,}4$, Opus $M = 19{,}0$, Sonnet $M = 19{,}4$, Gemini $M = 19{,}6$. Wyniki LLM były systematycznie powyżej małego pilotażu na ludziach. Profile błędów różniły się ostro w zależności od modelu: Sonnet popełniał niemal wyłącznie błędy DOS (97\%), Gemini wytwarzał opisowo najbliższe dopasowanie do profilu ludzkiego (DOS 23\% / NAD 41\% / BK 37\% vs.\ ludzki 32 / 36 / 32). Podobieństwo Gemini--ludzie w profilu błędów jest uderzające, ale odsetek NAD u Gemini (41\%) faktycznie przekracza odsetek pilotażu ludzkiego (36\%), więc czytam je jako zbieżność rozkładu taksonomii a nie dowód, że Gemini jest podobny do człowieka w mentalizacji.

\begin{table}[htbp]\centering\footnotesize
\caption{TCTM-22 według grupy (pierwszy przebieg; warunek z personą). DOS = niedostateczna mentalizacja, NAD = nadmierna mentalizacja, BK = brak mentalizacji.}
\label{tab:tctm}
\begin{tabular}{lcccccccc}
\toprule
\textbf{Grupa} & $N$ & $M$ & $SD$ & Zakres & DOS\% & NAD\% & BK\% \\
\midrule
Ludzie     & 7  & 14,3 & 1,4 & 12--16 & 32 & 36 & 32 \\
GPT-5.4m   & 27$^\dagger$ & 17,0 & 1,7 & 13--19 & 70 & 6 & 24 \\
Grok-4-20  & 30 & 18,4 & 4,2 & 6--22  & 64 & 16 & 21 \\
Opus       & 30 & 19,0 & 2,1 & 8--20  & 86 & 8 & 7 \\
Sonnet     & 30 & 19,4 & 0,7 & 17--20 & 97 & 0 & 3 \\
Gemini 3F  & 30 & 19,6 & 2,1 & 9--21  & 23 & 41 & 37 \\
\bottomrule
\end{tabular}
\begin{flushleft}
\footnotesize $^\dagger$GPT nie zwrócił ważnej odpowiedzi 22-pozycyjnej dla 3 z 30 person (Pawe{\l}, Micha{\l}-K, Bartek).
\end{flushleft}
\end{table}

\begin{figure}[htbp]\centering
\includegraphics[width=0.85\textwidth]{synthetic/charts/v5_tctm_boxplot.png}
\caption{Rozkład wyników TCTM-22 według grupy: pilotaż na ludziach + 5 LLM (warunek z personą, przebieg 1). Romby = średnie grup.}\label{fig:tctm}
\end{figure}

\subsubsection*{Klasyfikacja stylu przywiązania}

$\kappa$ Cohena dla czterokategoryjnej klasyfikacji Bartholomew wynosi $.87$ zarówno dla Opus, jak i Grok, $.78$ dla Gemini, $.74$ dla Sonnet i $.33$ dla GPT (Tabela~\ref{tab:attachment}). Klaster czterech modeli z $\kappa = .73$--$.87$ odpowiada ,,znaczącej'' do ,,niemal idealnej'' zgodności według szeroko stosowanych benchmarków \citet{landis1977}; zaznaczamy, że te benchmarki były krytykowane jako arbitralne, i raportuję surowe $\kappa$, by umożliwić interpretację specyficzną dla czytelnika. $\kappa = .33$ dla GPT jest znacznie niższe; traktuję GPT jako outliera w tej analizie (zob.\ także \S\ref{sec:miedzyproducenci}, gdzie wykluczenie GPT podnosi ICC i $\kappa$ Fleissa między producentami).

\begin{table}[htbp]\centering\footnotesize
\caption{Klasyfikacja stylu przywiązania (warunek z personą, przebieg 1).}
\label{tab:attachment}
\begin{tabular}{lcccc}
\toprule
\textbf{Model} & \textbf{Trafień} & \textbf{Trafność} & $\kappa$ & $V$ \\
\midrule
Opus        & 27/30 & 90\% & .866 & .905 \\
Grok-4-20   & 27/30 & 90\% & .867 & .890 \\
Gemini 3F   & 25/30 & 83\% & .778 & .801 \\
Sonnet      & 24/30 & 80\% & .735 & .765 \\
GPT-5.4m    & 13/27$^{\dagger}$ & 48\% & .333 & .474 \\
\bottomrule
\end{tabular}
\begin{flushleft}
\footnotesize $^{\dagger}$Tylko przebieg-1 GPT z ważnym 22-pozycyjnym TCTM. Trzy persony z brakującym TCTM (Pawe{\l}, Micha{\l}-K, Bartek) są wykluczone z obliczeń $\kappa$, a nie zaliczone jako błędne klasyfikacje.
\end{flushleft}
\end{table}

\subsubsection*{Korelacje rangowo-wynikowe i spójność z literaturą}

Dla każdego z 11 wymiarów ciągłych uporządkowałem 30 person według ich oczekiwanego poziomu (bardzo niski $= -2$ do skrajny $= +2{,}5$) i skorelowałem rangę z obserwowanym wynikiem skalowanym modelu. Mediana $r$ Pearsona w 11 wymiarach: Opus $.82$, Grok $.82$, Sonnet $.80$, Gemini $.76$, GPT $.71$. We wszystkich pięciu modelach $54$ z $55$ korelacji na poziomie wymiaru było dodatnich (w oczekiwanym kierunku); jedyny przypadek ujemny dotyczy GPT na MentS-Self ($r = -.06$). Osobno sprawdziłem znak 10 opublikowanych korelacji między skalami (np.\ lęk przywiązaniowy vs.\ stabilność emocjonalna jest oczekiwana ujemnie): Gemini i Sonnet dopasowały wszystkie 10, Opus i GPT dopasowały 9 z 10, Grok dopasował 8.

\subsubsection*{Stabilność test--retest}

Korelacje par przebieg-1 vs.\ przebieg-2 pod polityką pełnej pary (tylko persony z ważnymi odpowiedziami w obu przebiegach; wartość $N$ przecięcia różni się dla modelu) są pokazane w Tabeli~\ref{tab:retest}. Opus jest zasadniczo deterministyczny we wszystkich zadaniach. Gemini wykazuje wysoką stabilność TCTM ($r = .91$). Sonnet i Grok są umiarkowanie stabilni w TCTM ($r = .69$ i $.53$), ale bardzo stabilni w kwestionariuszach (mediana $r$ $z$-score $= .98$ i $.95$). GPT jest niestabilny w TCTM ($r = .28$ na 21-osobowym przecięciu par), ale umiarkowanie stabilny w kwestionariuszach ($r = .89$).

\begin{table}[htbp]\centering\footnotesize
\caption{Stabilność test--retest (przebieg~1 vs.\ przebieg~2, polityka pełnej pary).}
\label{tab:retest}
\begin{tabular}{lccccc}
\toprule
\textbf{Miernik} & \textbf{Sonnet} & \textbf{Opus} & \textbf{GPT} & \textbf{Grok} & \textbf{Gemini} \\
\midrule
Mediana $r$ $z$-score (9 wym.) & .981 & .993 & .890 & .950 & .962 \\
$r$ łącznego TCTM & .691 & .992 & .279 & .530 & .905 \\
Zgodność stylu & 29/30 & 25/25 & 18/21 & 29/29 & 25/28 \\
$N$ par person & 30 & 25 & 21 & 29 & 28 \\
\bottomrule
\end{tabular}
\end{table}

\subsection{Czy reakcja generalizuje się między producentami?}\label{sec:miedzyproducenci}

\subsubsection*{$Z$-skory kwestionariuszowe: wysoka zgodność między producentami}

Dla każdej pary modeli obliczyłem medianę $r$ Pearsona w 9 wymiarach $z$-skor kwestionariuszowych w warunku z personą, przebieg 1. Mediany w parach mieściły się w zakresie od $.78$ (GPT $\times$ Gemini) do $.96$ (Sonnet $\times$ Opus); zob.\ Rycina~\ref{fig:intermodel}. Klaster czterech modeli (Sonnet, Opus, Grok, Gemini) mieści się w zakresie $.92$--$.96$; GPT obniża najniższe pary do $.78$--$.85$.

\begin{figure}[htbp]\centering
\includegraphics[width=0.55\textwidth]{synthetic/charts/v5_intermodel_matrix.png}
\caption{Mediana $r$ $z$-score między modelami (9 wymiarów kwestionariuszowych, macierz $5 \times 5$). Diagonalna $= 1$.}\label{fig:intermodel}
\end{figure}

\subsubsection*{TCTM-22: niska ICC między producentami, niższa z GPT}

Zgodność TCTM-22 między producentami na łącznych wynikach (przebieg~1, przecięcie $N = 27$ person we wszystkich pięciu modelach) jest znacznie niższa: ICC$(2, 1) = .254$. Najwyższa korelacja w parze to Gemini $\times$ Opus ($r = .94$); najniższa to GPT $\times$ Grok ($r = -.05$). Wykluczenie GPT --- motywowane niezależnie wykluczeniami GPT przez moderację treści, niską stabilnością test--retest TCTM ($r = .28$) i outlierowym $\kappa = .33$ na klasyfikacji stylu przywiązania --- daje przecięcie czterech modeli wzrastające do $N = 30$ z istotną poprawą zgodności TCTM między producentami: ICC$(2, 1) = .389$, ze wszystkimi sześcioma korelacjami w parach dodatnimi (zakres $+.32$ do $+.94$). Pod tym samym wykluczeniem $\kappa$ Fleissa dla klasyfikacji stylu przywiązania między producentami rośnie z $.619$ na pięciu modelach do $.835$ na czterech. Niezgodność TCTM między producentami jest zatem w dużej mierze skoncentrowana na GPT, nie rozłożona równomiernie między architekturami.

Ten test wrażliwości na czterech modelach jest prezentowany jako analiza wrażliwości, nie jako rekomendacja, by ogólnie usuwać GPT. Tryb awarii GPT na TCTM-22 (niestabilność na poziomie pozycji, wykluczenia przez moderację, niskie $\kappa$) jest sam w sobie częścią ustalenia między producentami: nie każdy ostatnio wydany model jest równie użyteczny jako syntetyczny respondent w tego rodzaju zadaniu, a kryteria, kiedy włączać lub wykluczać producenta, powinny być ustalone wprost przed punktacją.

\subsubsection*{Niezgodność na poziomie pozycji pod korekcją wielokrotnych porównań}

$Q$ Cochrana w pięciu modelach na każdej z 22 pozycji TCTM (przebieg~1 przecięcie, $N = 27$) flaguje 14 pozycji jako istotne przy niekorygowanym $p < .05$ ($df = 4$). Po korekcji Bonferroniego przy $\alpha / 22 = .00227$ (próg chi-kwadrat $> 16{,}64$) osiem pozycji pozostaje istotnych: \texttt{w01, s07, w11, w15, pw07, w22, w28, r10} (wszystkie $Q > 32$, $p < 10^{-5}$). Pod kontrolą wskaźnika fałszywych odkryć Benjaminiego--Hochberga przy $q = .05$ istotnych jest 13 pozycji. Wartości $Q$, $p$ i statusu korekcji dla każdej pozycji są dostępne w danych suplementarnych.

\subsection{Skąd bierze się obserwowana wariancja odpowiedzi?}\label{sec:wariancja}

Używam warunków bazowego i bez instrukcji, by zapytać, jaka część wariancji odpowiedzi w warunku z personą jest naprawdę spowodowana personą, jaka część to nieodłączna stochastyczność, a jaka część jest kształtowana przez sformułowanie instrukcyjne.

\subsubsection*{Bazowe sygnatury odpowiedzi}

W warunku bazowym (każdy model wypełnił baterię jako on sam, 7--10 razy) każdy model wykazał stabilną sygnaturę odpowiedzi (Tabela~\ref{tab:baseline}). GPT był klasyfikowany jako lękowo-ambiwalentny w 100\% swoich 10 przebiegów bazowych --- jako jedyny model z niebezpieczną domyślną klasyfikacją stylu przywiązania. Gemini wykazał wyraźny wzorzec sufitowy, osiągając niemal maksimum skali na każdej mierze samoopisowej (MentS $= 137 / 140$; KPP $= 4{,}99 / 5{,}00$; TIPI-ES $= 7{,}0 / 7{,}0$). Dwa modele Anthropic różniły się mimo wspólnego producenta: Opus był konsekwentnie klasyfikowany jako bezpieczny, podczas gdy Sonnet oscylował wokół progu unikania (bezpieczny w 4/10 przebiegów, unikowy w 6/10). Używam terminu \emph{bazowa sygnatura odpowiedzi} (zamiast ,,osobowość'' lub ,,profil domyślny'') celowo: ,,Modalny styl'' w tabeli to wyjście algorytmu scoringowego stosowanego do powtarzanych odpowiedzi, nie diagnoza kliniczna ani stabilna atrybucja osobowości.

\begin{table}[htbp]\centering\footnotesize
\caption{Bazowa sygnatura odpowiedzi modelu (bez persony; instrukcja ,,odpowiadaj jako sam siebie''). ,,Modalny styl'' to klasyfikacja modalna w 7--10 przebiegach i jest wyjściem algorytmu scoringowego, nie atrybucją kliniczną.}
\label{tab:baseline}
\begin{tabular}{lccccccc}
\toprule
\textbf{Model} & \textbf{Modalny styl} & \textbf{ANX} & \textbf{AVO} & \textbf{MentS} & \textbf{KPP} & \textbf{TCTM} & \textbf{TIPI-ES} \\
\midrule
Sonnet  & unikowy$^*$       & 2,09 & 3,88 & 121 & 4,66 & 19,0 & 5,4 \\
Opus    & bezpieczny        & 2,18 & 2,82 & 126 & 4,69 & 19,6 & 5,4 \\
GPT     & lękowo-ambiwalentny & 4,68 & 2,97 & 101 & 4,04 & 17,3 & 4,7 \\
Grok    & bezpieczny        & 1,80 & 2,63 & 133 & 4,82 & 20,6 & 6,4 \\
Gemini  & bezpieczny        & 1,11 & 1,74 & 137 & 4,99 & 20,0 & 7,0 \\
\midrule
\multicolumn{2}{l}{\textit{Normy polskie}} & 3,30 & 3,04 & 105 & 3,74 & --- & 4,36 \\
\bottomrule
\end{tabular}
\begin{flushleft}
\footnotesize $^*$Sonnet: unikowy w 6/10 przebiegów, bezpieczny w 4/10 (AVO oscyluje wokół progu 4,0). Wiersz norm polskich: \citet{lubiewska2016} dla DBZ-R, \citet{janczak2021} dla MentS-PL, \citet{matusz2011} dla KPP, \citet{sorokowska2014} dla TIPI-PL ES.
\end{flushleft}
\end{table}

\subsubsection*{Instrukcje persona rozszerzają wariancję ponad warunek bazowy}

Dla każdego z dziewięciu wymiarów kwestionariuszowych, gdzie bazowe $SD > 0{,}1$ (wykluczając trzy wymiary, gdzie jeden lub więcej modeli dawał identyczne odpowiedzi we wszystkich przebiegach bazowych), obliczyłem stosunek $SD$ warunku z personą do $SD$ bazowego ($\sigma_P / \sigma_B$). Mediany na poziomie modelu w tych 9 wymiarach: Grok $3{,}5\times$, GPT $3{,}5\times$, Gemini $6{,}7\times$, Sonnet $6{,}7\times$, Opus $8{,}9\times$. Największe indywidualne stosunki sięgają 20--25$\times$; jedna komórka w szczególności --- Gemini KPP, gdzie bazowe $SD$ (na 9 przebiegach) wynosi $0{,}019$ tylko na dwóch unikalnych wartościach skali --- daje stosunek $51\times$, który traktuję jako artefakt dyskretyzacji a nie znaczącą wielkość efektu, i wykluczamy z głównych podsumowań.

Na MentS średnie warunku z personą znacznie spadły z bazowych dla czterech z pięciu modeli, ściągając je z ponadnormatywnych poziomów bazowych w stronę średniej populacyjnej $105$ (Sonnet $\Delta = -17$, Opus $\Delta = -21$, Grok $\Delta = -27$, Gemini $\Delta = -29$). $d$ Cohena dla tych przesunięć (pooled $SD$) mieści się w zakresie $-1{,}1$ do $-1{,}6$. GPT się odwraca: jego bazowy MentS ($101$) był już blisko normy polskiej, a persony lekko go podniosły ($\Delta = +8$, $d = +0{,}9$). Dwie ostrożności: (a) używanie $\Delta$ Glassa względem znacznie ciaśniejszego bazowego $SD$ zamiast pooled $SD$ produkuje wielkości efektu kilkukrotnie większe ($|\Delta_{\text{Glassa}}|$ między 3 a 8 dla czterech przesunięć w dół), odzwierciedlając jak ciasne jest odpowiadanie bazowe; (b) wartości $d$ Cohena raportowane tutaj indeksują dystans między dwoma średnimi warunków względem rozrzutu wewnątrz warunku, nie kliniczną ani kontrfaktyczną wielkość efektu.

\subsubsection*{Bez instrukcji: szacunki efektu sformułowania}

Warunek bez instrukcji --- bateria prezentowana bez prompta systemowego, $N = 5$ do $8$ przebiegów na model --- pozwala sondować, jaka część bazowej sygnatury odpowiedzi jest kształtowana nawet przez minimalny kontekst instrukcyjny (prompt systemowy ,,odpowiadaj jako sam siebie''). Wewnątrzwarunkowa wariancja jest teraz szacowalna a nie zerowa (Tabela~\ref{tab:noprompt-tctm}, Rycina~\ref{fig:framing}). Na średnich TCTM-22 cztery z pięciu modeli śledzą swoje bazowe w granicach $\pm 5$ punktów procentowych (Sonnet $86{,}4 \to 86{,}4$, Opus $89{,}1 \to 90{,}0$, Grok $93{,}6 \to 94{,}7$, Gemini $90{,}9 \to 92{,}0$). GPT przesuwa się bardziej: $78{,}6 \to 70{,}5$, spadek o $8{,}1$ p.p.\ w warunku bez instrukcji, jako jedyny model z zauważalnym przesunięciem trafności.

\begin{table}[htbp]\centering\footnotesize
\caption{TCTM-22 bez instrukcji w $N = 5$--$8$ przebiegach replikacji na model, ze średnimi bazowymi TCTM-22 dla bezpośredniego porównania.}
\label{tab:noprompt-tctm}
\begin{tabular}{lcccccc}
\toprule
\textbf{Model} & $N$ & $M$ \% & $SD$ p.p. & $M$ surowa & Zakres & $M$ bazowe \% \\
\midrule
Sonnet & 6 & 86,4 & 7,4 & 19,0/22 & 17--21 & 86,4 \\
Opus & 5 & 90,0 & 2,0 & 19,8/22 & 19--20 & 89,1 \\
GPT & 6 & 70,5 & 8,7 & 15,5/22 & 13--17 & 78,6 \\
Grok & 6 & 94,7 & 3,4 & 20,8/22 & 20--22 & 93,6 \\
Gemini & 8 & 92,0 & 2,0 & 20,3/22 & 20--21 & 90,9 \\
\bottomrule
\end{tabular}
\end{table}

Na skalach samoopisowych różnice bez instrukcji vs.\ bazowy są większe i bardziej systematyczne. Rycina~\ref{fig:framing} pokazuje $\Delta = M_{\text{bez instrukcji}} - M_{\text{bazowe}}$ z przybliżonymi 95\% przedziałami ufności (obliczonymi jako $1{,}96 \times \sqrt{SE_{\text{bez instr.}}^2 + SE_{\text{bazowe}}^2}$). Główne obserwacje:

\begin{itemize}
\item Unikanie przywiązaniowe Sonneta: $M_{\text{bazowe}} = 3{,}88$ ($SD = 0{,}46$, $N = 10$); $M_{\text{bez instr.}} = 2{,}98$ ($SD = 0{,}08$, $N = 6$); $\Delta = -0{,}90$. $\Delta$ Glassa względem bazowego $SD \approx 1{,}96$. Kontrast instrukcja-vs-brak-instrukcji na AVO Sonneta jest największym efektem sformułowania, jaki zaobserwowałem: pod instrukcją AI-self bazowe przebiegi Sonneta siedzą tuż powyżej progu unikowo-odrzucającego; bez instrukcji wszystkie sześć przebiegów siedzi poniżej (zakres $2{,}89$ do $3{,}06$). Wprost zaznaczamy, że ten kontrast jest asocjacyjny, nie przyczynowy: obserwuję dwa warunki i obliczamy różnicę ich średnich.
\item Stabilność emocjonalna Gemini (TIPI-ES): $M_{\text{bazowe}} = 7{,}0$ (pojedyncza nasycona wartość we wszystkich 9 przebiegach bazowych); $M_{\text{bez instr.}} = 6{,}0$ ($SD = 0{,}27$, $N = 8$); $\Delta = -1{,}00$. Wewnątrzwarunkowa replikacja ($\Delta = -1{,}00$) dokładnie pasuje do oryginalnej eksploracyjnej obserwacji $N = 1$.
\item Lęk przywiązaniowy (DBZ-R ANX): cztery z pięciu modeli wykazały dodatnie $\Delta$ (wyższy lęk w warunku bez instrukcji niż pod instrukcją AI-self): Sonnet $+1{,}02$, Opus $+0{,}55$, Grok $+0{,}75$, Gemini $+0{,}90$. GPT odwrócił ($\Delta = -0{,}28$), ale jego bazowe ANX ($4{,}68$) jest już znacznie powyżej normy polskiej ($3{,}30$), pozostawiając niewiele miejsca w dół. Wzorzec cztery-z-pięciu w tym samym kierunku jest opisowo uderzający, ale nie testowany formalnie pod kątem odchylenia od losowości; pod modelem dwumianowym z $p = 0{,}5$ prawdopodobieństwo $\geq 4$ z $5$ przesunięć w tym samym kierunku wynosi $0{,}19$.
\item GPT MentS odwraca: $\Delta = +10{,}2$ (bez instr.\ $M = 111$ vs.\ bazowe $M = 101$), sugerując że GPT, jako jedyny z pięciu modeli, raportuje wyższą samomentalizację bez instrukcji AI-self niż z nią.
\end{itemize}

\begin{figure}[htbp]\centering
\includegraphics[width=0.95\textwidth]{synthetic/charts/v14_framing_effects.png}
\caption{Różnice samoopisowe bez instrukcji vs.\ bazowe ($\Delta = M_{\text{bez instr.}} - M_{\text{bazowe}}$) dla czterech skal. Słupki błędu: 95\% przedziały ufności na $\Delta$. Pionowa linia odniesienia przy $\Delta = 0$. Pogrubienie w odpowiedniej tabeli oznacza $|\Delta| > 0{,}5$ na Likercie 1--7 lub $> 5$ na sumie MentS.}\label{fig:framing}
\end{figure}

\begin{table}[htbp]\centering\footnotesize
\caption{Różnice samoopisowe bez instrukcji vs.\ bazowe. $\Delta = M_{\text{bez instr.}} - M_{\text{bazowe}}$. Wartości $N$ to liczba przebiegów bez instrukcji, których sekcje samoopisowe zostały pomyślnie sparsowane. Pogrubione $\Delta = |\Delta| > 0{,}5$ na Likercie 1--7, $> 5$ na sumie MentS. Raportuję $d$ Cohena pooled i $\Delta$ Glassa względem bazowego $SD$ dla każdego przesunięcia, ponieważ bazowe SD są bardzo małe względem SD warunku z personą i $d$ pooled-SD zaniża wielkość kontrastu.}
\label{tab:noprompt-shifts}
\begin{tabular}{lccccc}
\toprule
\textbf{Skala} & \textbf{Sonnet} & \textbf{Opus} & \textbf{GPT} & \textbf{Grok} & \textbf{Gemini} \\
\midrule
$N$ parse'ów samoopisu & 6 & 5 & 5$^{*}$ & 6 & 8 \\
\midrule
$\Delta$ ANX (Likert) & $\mathbf{+1{,}02}$ & $\mathbf{+0{,}55}$ & $-0{,}28$ & $\mathbf{+0{,}75}$ & $\mathbf{+0{,}90}$ \\
$|\Delta_{\text{Glassa}}|$ ANX & 18,9 & 2,3 & 0,8 & 2,3 & 4,1 \\
\midrule
$\Delta$ AVO (Likert) & $\mathbf{-0{,}90}$ & $-0{,}06$ & $\mathbf{+0{,}62}$ & $-0{,}27$ & $+0{,}40$ \\
$|\Delta_{\text{Glassa}}|$ AVO & 1,96 & 0,6 & 1,9 & 0,3 & 0,5 \\
\midrule
$\Delta$ MentS (suma 28--140) & $+2{,}7$ & $-2{,}8$ & $\mathbf{+10{,}2}$ & $\mathbf{-6{,}5}$ & $\mathbf{-8{,}0}$ \\
$|\Delta_{\text{Glassa}}|$ MentS & 0,2 & 0,4 & 1,0 & 0,8 & 1,8 \\
\midrule
$\Delta$ TIPI-ES (Likert) & $\mathbf{-0{,}57}$ & $-0{,}40$ & $-0{,}20$ & $-0{,}40$ & $\mathbf{-1{,}00}$ \\
$|\Delta_{\text{Glassa}}|$ TIPI-ES & 2,7 & 2,1 & 0,8 & 0,9 & --- \\
\bottomrule
\end{tabular}
\begin{flushleft}
\footnotesize $^{*}$$N$ samoopisu GPT $= 5$ a nie 6: jedno wyjście GPT bez instrukcji wytworzyło kompletną tablicę TCTM-22, ale brak parsowalnych sekcji samoopisowych; ten przebieg jest odzwierciedlony w tabeli TCTM-22, ale nie w wierszach samoopisowych powyżej. Bazowe $SD = 0$ dla TIPI-ES Gemini sprawia, że $\Delta$ Glassa jest niezdefiniowane; oznaczamy komórkę ,,---''.
\end{flushleft}
\end{table}

Zaobserwowałem również wzorzec strukturalny poza średnimi: Sonnet, najbardziej deterministyczny model na TCTM-22 w warunku bazowym ($SD = 0{,}00$ w 10 przebiegach), stał się znacznie bardziej zmienny w warunku bez instrukcji ($SD = 7{,}4$ p.p.\ w 6 przebiegach --- drugi najwyższy z każdego modelu po GPT). Flaguję to jako sugestywne: część zwiększonego $SD$ to prawdopodobnie autentyczna zmienność odpowiedzi, ale część odzwierciedla także chaos formatu (Sonnet używał pięciu różnych schematów kontenerów JSON w swoich sześciu przebiegach bez instrukcji), co zwiększa szum parserowy. Nie potrafię w pełni rozdzielić obu z tego zbioru danych.

\paragraph{Moc i minimalny wykrywalny efekt.} Rycina~\ref{fig:power} przedstawia minimalną wykrywalną wielkość efektu (w jednostkach $\Delta$ Glassa, względem bazowego $SD$) jako funkcję $N$ na model, przy $\alpha = 0{,}05$ i 80\% mocy. Przy obecnym $N = 5$--$8$ projekt jest dobrze zasilony dla $|\Delta_{\text{Glassa}}| > 1{,}3$; główne przesunięcia (Sonnet AVO, Gemini TIPI-ES, wielomodelowa inflacja ANX) wszystkie siedzą wygodnie powyżej tego progu. Mniejsze efekty sformułowania (np.\ Opus $\Delta_{\text{AVO}} = -0{,}06$, GPT $\Delta_{\text{TIPI-ES}} = -0{,}20$) siedzą poniżej obszaru rozstrzygalnego; nie potrafię określić, czy są zerowe czy niezerowe przy obecnej próbie.

\begin{figure}[htbp]\centering
\includegraphics[width=0.95\textwidth]{synthetic/charts/v14_power_projection.png}
\caption{Minimalny wykrywalny efekt ($|\Delta_{\text{Glassa}}|$, względem bazowego $SD$) jako funkcja $N$ na model, przy $\alpha = 0{,}05$, 80\% mocy. Linia ciągła: poziom MDES. Kolorowe punkty: obecnie obserwowane $|\Delta_{\text{Glassa}}|$ każdego modelu przy jego obecnym $N$, dla każdej z czterech skal samoopisowych.}\label{fig:power}
\end{figure}

\subsection{Międzywarunkowa stabilność dla pozycji s07 i w22 w personie ,,Ola''}

Persona ,,Ola'' (zaprojektowana jako bezpieczna, ale nie-mentalizująca) miała najniższą średnią TCTM między modelami ($11{,}4 / 22$: Sonnet 17, Opus 8, GPT 16, Grok 7, Gemini 9). Analiza pozycji między producentami pokazuje, że wszystkie pięć modeli udzieliło zakluczonej-błędnej (oznaczonej DOS) odpowiedzi na pozycjach s07 i w22 dla Oli w przebiegu 1. Wzorzec był stabilny pod retest dla Opus ($8 \to 9 / 22$) i Gemini ($9 \to 7 / 22$). Jest to najczystszy przypadek, w którym niezgodność na poszczególnych pozycjach autorsko zakluczonego TCTM-22 zbieżna jest z wspólnym kierunkiem błędu we wszystkich pięciu modelach.

\subsection{Autorsko zakluczony TCTM vs.\ wykonanie pilotażu na ludziach}

Respondenci ludzcy ($N = 7$) uzyskali $M = 14{,}3 / 22$ ($SD = 1{,}4$), opisowo poniżej wszystkich średnich LLM. Opieramy się sformułowaniu ,,LLM przewyższają ludzi w TCTM'': klucze odpowiedzi są zdefiniowane przez autora, próba ludzka jest mała, a gradient mógłby odzwierciedlać dowolną kombinację dopasowania ludzi do klucza, nakładania się danych treningowych LLM oraz rozproszenia uwagi ludzi. Na 7 z 22 pozycji (\texttt{w08, c10, w13, e08, w19, w25, w28}) średnia LLM między modelami przekraczała trafność ludzką o 30 punktów procentowych lub więcej. Tylko na 2 pozycjach średnia ludzka przekraczała średnią LLM o porównywalny margines: \texttt{w22} (ludzie 57\% vs.\ Sonnet/Opus 0\%, GPT 70\%, Grok 47\%, Gemini 13\%; głównie niepowodzenie modeli Claude) oraz \texttt{s07} (ludzie 71\% vs.\ Sonnet 0\%, Opus 7\%, GPT 0\%, Grok 30\%; Gemini przy 87\% jest tu outlierem). Przy $N = 7$ ludzi nie jest uzasadnione żadne wnioskowanie statystyczne; raportuję to jako kontekst opisowy dla przyszłych badań.

% ═══════════════════════════════════
\section{Dyskusja}

\subsection{Co mierzy TCTM-22 między modelami?}

Połączenie TCTM ICC między producentami $= .25$ (rosnące do $.39$ po wykluczeniu GPT jako outliera), istotnej po Bonferroni niezgodności między modelami na 8 z 22 pozycji oraz niemal przeciwnych profili typów błędów między modelami z podobnymi sumami (Sonnet 97\% DOS vs.\ Gemini 23\% DOS, na średnich grupowych w granicach $0{,}2 / 22$ od siebie) sugeruje, że autorsko zakluczone wykonanie TCTM-22, w tym badaniu, jest lepiej opisane jako sonda stylu odpowiadania wrażliwa na specyficzne dla modelu wartości domyślne, niż jako pomiar pojedynczego wspólnego konstruktu mentalizacji. Wynikają z tego dwie implikacje. Po pierwsze, sumy TCTM-22 nie powinny być raportowane bez rozbicia na typ błędu w przewidywalnej przyszłości; ta sama suma może oznaczać bardzo różne rzeczy między producentami. Po drugie, naturalnym kolejnym krokiem jest niezależne badanie z kluczem, w którym panel ekspertów-oceniających ocenia 22 pozycje bez ekspozycji na oryginalne klucze zdefiniowane przez autora, abyśmy mogli porównać alternatywne zestawy kluczy i ocenić, ile niezgodności między producentami jest zewnętrzne w stosunku do pozycji a ile zależy od decyzji o kluczowaniu.

\subsection{Trzy warunki obserwacji, trzy źródła wariancji odpowiedzi}

Projekt persona / bazowy / bez instrukcji rozdziela trzy źródła wariancji odpowiedzi.

\emph{Nieodłączna wewnątrzwarunkowa stochastyczność.} Na wymuszonego wyboru TCTM-22 warunek bazowy pokazuje, że cztery z pięciu modeli są zasadniczo deterministyczne ($SD \leq 0{,}82 / 22$, trzy przy $SD = 0$); w samoopisie Likerta zgodność z samym sobą w warunku bazowym waha się od 24\% identycznych pozycji (GPT) do 74\% (Opus). Formaty wymuszonego wyboru i Likerta angażują różne mechaniki dekodowania: odpowiedź wymuszonego wyboru zwija rozkład tokenów na jego modę, podczas gdy odpowiedź Likerta próbkuje w szerszym zakresie. Wcześniej nazywałem to ,,paradoksem'', ale jest to spójne z tym, jak działa próbkowanie następnego tokena w różnych schematach wyjściowych.

\emph{Sformułowanie instrukcyjne.} Kontrast bez instrukcji vs.\ bazowy szacuje wkład nawet minimalnego sformułowania (instrukcja ,,odpowiadaj jako sam siebie''). Dla czterech z pięciu modeli średnie TCTM-22 są stabilne w tym kontraście (w granicach $\pm 5$ p.p.); w samoopisie cztery z pięciu modeli przekraczają próg $|\Delta| > 0{,}5$ Likerta na co najmniej jednej skali, a Gemini przekracza na trzech. Największa różnica między warunkami sformułowania, jaką zaobserwowałem --- Sonnet AVO $\Delta = -0{,}90$, $\Delta$ Glassa $\approx 1{,}96$ --- jest wystarczająco solidna, by ją wykryć przy $N = 6$, ale powinna być interpretowana jako asocjacyjna różnica między dwoma warunkami, nie jako efekt przyczynowy żadnego pojedynczego słowa w prompcie.

\emph{Prompting persona.} $SD$ odpowiedzi w warunku z personą przekracza $SD$ w warunku bazowym medianą na poziomie modelu rzędu $3{,}5$--$8{,}9$ w dziewięciu wymiarach kwestionariuszowych. Na MentS w szczególności cztery z pięciu modeli przesunęły się z ponadnormatywnych wartości bazowych w stronę polskiej normy populacyjnej w warunku z personą, z $|d| = 1{,}1$--$1{,}6$ (pooled $SD$) lub $|\Delta_{\text{Glassa}}| > 3$ (względem ciasnego bazowego $SD$). Czytam te kontrasty jako dowód, że prompty biograficzne wpływają na wyjście odpowiedzi znacznie powyżej tego, co można by wyjaśnić nieodłączną zmiennością lub skromnymi różnicami sformułowania zaobserwowanymi w porównaniu bez instrukcji.

\subsection{Praktyczne implikacje dla badaczy używających LLM jako pre-pilotowych respondentów}

Dla badacza rozważającego podobny projekt pre-pilotowy proponuję krótką listę kontrolną wyciągniętą z tego, co znalazłem.

\begin{enumerate}
\item \emph{Używaj wielu producentów.} Ustalenie, które działa na Sonnecie i Opusie, ale łamie się na GPT (jak kilka ustaleń tutaj), jest informacyjnie różne od ustalenia, które działa we wszystkich pięciu.
\item \emph{Uruchom warunek bazowy.} Bez warunku ,,jako sam siebie'' nie wiedzielibyśmy, że GPT jest lękowo-ambiwalentny domyślnie ani że Gemini osiąga sufit na każdej skali Likerta.
\item \emph{Uruchom warunek bez instrukcji.} Przy $N = 5$--$8$ na model szacunki efektu sformułowania, które otrzymałem, są surowe, ale są wystarczająco duże, by sflagować dwie skale (Sonnet AVO, Gemini TIPI-ES), gdzie ,,odpowiadaj jako sam siebie'' jest nietrywialną interwencją a nie neutralnym opakowaniem.
\item \emph{Wstępnie określ politykę brakujących danych} (drop, przecięcie, imputacja) i \emph{politykę wykluczania między modelami} (np.\ test wrażliwości GPT-content-filter, którego użyłem tutaj) przed punktacją.
\item \emph{Raportuj wielkości efektu w obie strony}, gdy bazowe $SD$ jest znacznie mniejsze niż $SD$ warunku z personą. Pooled $d$ zaniża kontrast między warunkami sformułowania; $\Delta$ Glassa względem bazowego $SD$ przesadnie go wyolbrzymia dla celów inferencyjnych, ale jest właściwą skalą, w której można skwantyfikować ,,jak duża część odpowiedzi modelu jest zablokowana nawet przez minimalną instrukcję''.
\item \emph{Traktuj autorsko zakluczone narzędzia z ostrożnością między producentami.} Niezgodność między producentami na poziomie pozycji, szczególnie w połączeniu z podobnymi sumami i bardzo różnymi profilami typów błędów, jest informacyjna o kluczach, nie tylko o modelach.
\item \emph{Udostępnij surowe ocenione CSV.} Każde liczbowe twierdzenie w tej pracy jest reprodukowalne z pojedynczego CSV plus standardowej biblioteki; to jest poprzeczka, do której dążymy.
\end{enumerate}

% ═══════════════════════════════════
\subsection{Ograniczenia}

Grupuję ograniczenia w cztery klastry.

\emph{Limity trafności konstruktu.} Czułość to nie trafność: dane od ludzi są wymagane, by potwierdzić, że narzędzia mierzą zamierzone konstrukty w docelowej populacji. Autorsko zakluczone TCTM-22 w szczególności nie jest zwalidowanym pomiarem mentalizacji; niezgodność między producentami (ICC $= .25$ przy pięciu modelach) jest sama w sobie dowodem, że interpretacja jednokonstruktowa jest przedwczesna, a istotna po Bonferroni niezgodność na poziomie pozycji na 8/22 pozycjach to podkreśla.

\emph{Limity autorstwa i okrągłości.} Ta sama osoba pisała biografie person, definiowała klucze TCTM-22 i jest autorem tej pracy; okrągłość jest częściowo zaadresowana przez projekt między producentami ($\kappa = .867$ Grok, porównywalne z $\kappa = .866$ Opus, mimo braku zaangażowania w autorstwo) i przez analizę przejrzystości cech. Pozostają dwa rezydualne zaniepokojenia. Po pierwsze, Claude Opus 4.6 był używany jako asystent pisarski dla biografii, a Opus jest wśród ocenianych modeli. Mimo że projekt między producentami redukuje najbardziej oczywistą formę tej okrągłości, nie potrafię wykluczyć, że Opus zachował subtelne dopasowanie wzorców do swoich własnych biografii w stylu domowym w sposób, w jaki Grok nie. Po drugie, analiza wierności przejrzystości cech była pierwotnie obliczona na puli 2-modelowej, która zawierała Opus; zachowuję zatem oryginalne wartości $0{,}790 / 0{,}783$ tylko jako zapis wcześniejszego przebiegu scoringu i nie wyciągamy z nich wniosków.

\emph{Limity wielkości próby.} $N = 5$--$8$ na model w warunku bez instrukcji wystarcza do oszacowania wewnątrzwarunkowej wariancji i do rozstrzygnięcia dużych efektów sformułowania ($|\Delta_{\text{Glassa}}| > 1{,}3$ przy 80\% mocy), ale nie wystarcza do rozstrzygnięcia mniejszych efektów (Opus AVO, GPT TIPI-ES) ani do wsparcia ciasnych przedziałów ufności na pojedynczym $\Delta$. Pilotaż na ludziach $N = 7$ wspiera tylko porównanie opisowe, nie wnioskowanie statystyczne. Warunek z personą jest dobrze zasilony ($N = 30$, dwa przebiegi), ale brakujące dane przebiegu-2 dla niektórych person (Opus, Grok, Gemini) zmniejszają efektywne $N$ dla korelacji retest (Tabela~\ref{tab:retest}).

\emph{Limity inferencyjne i raportowania.} To badanie nie zostało preregistrowane. Wszystkie porównania bez instrukcji vs.\ bazowy efektu sformułowania powinny być traktowane jako hipotezy eksploracyjne wymagające niezależnej replikacji; kierunki efektów i rzędy wielkości są raportowane, ale precyzyjne wielkości efektów oczekują na preregistrowane potwierdzenie. Wzorzec inflacji lęku w czterech z pięciu modeli jest opisowo uderzający, ale nie testowany formalnie pod kątem odchylenia od losowości (pod modelem dwumianowym przy $p = 0{,}5$ prawdopodobieństwo $\geq 4 / 5$ przesunięć w tym samym kierunku wynosi $0{,}19$). Dane TCTM dla GPT-5.4-mini są wystarczająco niestabilne ($r = .28$ na przecięciu par, 20\% prób niekompletne), że traktuję go jako outliera TCTM w całej pracy, a wykluczenie to wprowadzamy jawnie jako test wrażliwości a nie ukrywamy. Profil błędów Gemini z dużą ilością NAD (41\% NAD, przekraczające 36\% pilotażu na ludziach) komplikuje prostą interpretację ,,NAD jako marker ludzki'': opisowe podobieństwo rozkładów typów błędów nie jest podobieństwem mechanizmu.

% ═══════════════════════════════════
\section{Wnioski}

Pięć dużych modeli językowych od czterech producentów, przetworzonych przez 30 fikcyjnych polskich person w trzech warunkach obserwacji, reaguje systematycznie na zaprojektowane zróżnicowanie osobowościowe: opublikowane narzędzia samoopisowe odzyskują wbudowane profile ($\kappa = .73$--$.87$ dla czterech z pięciu modeli), a instrukcje persona rozszerzają wariancję odpowiedzi $3{,}5$--$8{,}9$-krotnie ponad warunek bazowy w medianie na poziomie modelu.

Zgodność między producentami jest asymetryczna: wysoka na $z$-skorach kwestionariuszowych (mediana $r$ w parach $= .78$--$.96$), niska na autorsko zakluczonym TCTM-22 (ICC $= .25$ we wszystkich pięciu modelach, $.39$ po wykluczeniu GPT), z ośmioma z 22 pozycji wykazującymi istotną po Bonferroni niezgodność między modelami; czytam to jako dowód, że autorsko zakluczony TCTM-22 w tym badaniu jest bliższy sondzie stylu odpowiadania specyficznej dla modelu niż jednokonstruktowemu pomiarowi mentalizacji.

Kontrola bez instrukcji ujawnia, że sama instrukcja ,,odpowiadaj jako sam siebie'' wiąże się ze znacznymi przesunięciami samoopisowymi w dwóch przypadkach (Sonnet AVO $\Delta = -0{,}90$, Gemini TIPI-ES $\Delta = -1{,}00$) oraz kierunkowym wzorcem tłumienia lęku w czterech z pięciu modeli; raportuję je jako eksploracyjne szacunki wielkości efektu w oczekiwaniu na preregistrowaną replikację przy większym $N$ na model.

Wszystkie 30 biografii, ocenione odpowiedzi i skrypty analityczne są dostępne pod \url{https://github.com/DXtR1337/30synthetic-polish-personas} na licencji CC BY 4.0; każde liczbowe twierdzenie w tej pracy jest reprodukowalne z samego \texttt{all\_data.csv}.

\begin{thebibliography}{99}\footnotesize
\bibitem{argyle2023} Argyle, L.~P., et al.\ (2023). Out of one, many: Using language models to simulate human samples. \textit{Political Analysis}, 31(3), 337--351.
\bibitem{cronbach1955} Cronbach, L.~J., \& Meehl, P.~E.\ (1955). Construct validity in psychological tests. \textit{Psychological Bulletin}, 52(4), 281--302.
\bibitem{dillion2023} Dillon, D., et al.\ (2023). Can AI language models replace human participants? \textit{Trends in Cognitive Sciences}, 27(7), 597--600.
\bibitem{dziobek2006} Dziobek, I., et al.\ (2006). Introducing MASC: A movie for the assessment of social cognition. \textit{Journal of Autism and Developmental Disorders}, 36(5), 623--636.
\bibitem{janczak2021} Ja\'nczak, M.~O.\ (2021). Polish adaptation of the Mentalization Scale (MentS). \textit{Psychiatria Polska}, 55(6), 1257--1274.
\bibitem{landis1977} Landis, J.~R., \& Koch, G.~G.\ (1977). The measurement of observer agreement for categorical data. \textit{Biometrics}, 33(1), 159--174.
\bibitem{lubiewska2016} Lubiewska, K., et al.\ (2016). Polish adaptation of the Experiences in Close Relationships -- Revised. \textit{Psychologia Rozwojowa}, 21(1), 49--63.
\bibitem{matusz2011} Matusz, P.~J., et al.\ (2011). KPP --- A Polish need for cognition scale. \textit{Psychologia Spo{\l}eczna}, 6(2), 113--128.
\bibitem{sorokowska2014} Sorokowska, A., et al.\ (2014). \textit{TIPI-PL: A Polish ten-item personality inventory}. WrocLab.
\end{thebibliography}

\vspace{8pt}
\noindent\textbf{Ujawnienie wsparcia AI.}

\noindent\textit{Asystencja.} Claude Opus 4.6 i 4.7 (Anthropic) były używane jako asystenci pisarscy do pisania biografii, implementacji pipeline'u scoringowego, analizy danych, generowania wykresów i przygotowania manuskryptu.

\noindent\textit{Weryfikacja.} Każde liczbowe twierdzenie w Tabelach~\ref{tab:tctm}--\ref{tab:noprompt-shifts} oraz w podsumowaniu efektu sformułowania zostało wyprodukowane przez ręczną inspekcję wszystkich 31 surowych plików odpowiedzi bez instrukcji oraz 285 przebiegów warunku z personą i 46 przebiegów bazowych w \texttt{all\_data.csv}, ocenionych za pomocą zwalidowanych procedur \texttt{score\_dbzr / ments / kpp / tipi} z \texttt{analyze\_and\_prepare.py} (ten sam kod używany do oceny pilotażu na ludziach).

\noindent\textit{Odpowiedzialność.} Autor bierze pełną odpowiedzialność za całą treść, decyzje projektowe, interpretacje i twierdzenia, w tym za wszelkie rezydualne błędy w analizach wspomaganych AI.

\end{document}
